\scanpage{14}
\mspage{7}
% ll. 1–5
pleinement fidèle. Je vais par contre exiger que (pour tt $\pi$) la catégorie $\widetilde\pi\backslash \unsure{E|B}$ soit stable par $\varprojlim$ finies \inserted{non vides}, et que le foncteur $(*)$ soit exact à gauche.
% ll. 5–15
Prenant pour $\pi$ le foncteur constant de valeur $e$, on trouve que $\widetilde\pi\backslash B\simeq\pi\backslash B\simeq B$, on veut donc que $B$ soit stable par $\varprojlim$ finies \inserted{\uline{non vides}}, i.e. $A$ stable par $\unsure{\varinjlim|\underline{\lim}|\varprojlim}$ finies \uline{non vides}. Pour $B$, cela revient à la stabilité a) par produit binaire, et b) par produit fibré binaire (ou, au choix, par noyaux \unsure{des|de} doubles-flèches). Montrons qu'alors $\widetilde\pi\backslash B$ est stable par les \unsure{$\overline{\rm ms}$|mêmes} opérations.
% ll. 16–20
En fait, il faudra faire une hypothèse sur $\pi$, assurant que $\deleted{\unsure{f|c}}C=\widetilde\pi\backslash B$ est non seulement cofibré sur $B$, mais aussi \uline{fibré}. Cette hypothèse, c'est que $\pi$ transforme flèches de $B$ en
% archive sheet number: 13
